% Lapisan solusi O001 -- Bab 10 (Distribusi)
% 11 solusi asli terpisah; bukan tulisan John M. Erdman.
% Provenans model: OpenAI Codex gpt-5.6-sol, Ultra, atas arahan pengguna.
% Lisensi komponen solusi: CC BY-SA 4.0.
% Tidak menyiratkan dukungan John M. Erdman atau Portland State University.

\markboth{SOLUSI O001 UNTUK BAB 10: DISTRIBUSI}{SOLUSI O001 UNTUK BAB 10: DISTRIBUSI}
\section*{Solusi O001 untuk Bab 10: Distribusi}
\addcontentsline{toc}{section}{Solusi O001 untuk Bab 10: Distribusi}
\begin{center}
\small
Komponen solusi asli terpisah, CC BY-SA 4.0. Ditulis dengan bantuan
OpenAI Codex gpt-5.6-sol, Ultra, atas arahan pengguna. Pernyataan latihan
berasal dari edisi terjemahan karya John M. Erdman; solusi ini bukan tulisan
beliau dan tidak menyiratkan dukungan beliau atau Portland State University.
\end{center}

% O001-SOLUTION-ID: O001-FAOA-2015-CH10-EX-001-SOLUTION
% SOURCE-EXERCISE-ID: FAOA-2015-CH10-NODE-0034
% STATEMENT-TARGET-SHA256: 993392326124fdc4d83c44ae53cdb333a2ca2c9911da027d164621951aafc62a
\begin{o001solution}
  {O001-FAOA-2015-CH10-EX-001-SOLUTION}
  {FAOA-2015-CH10-NODE-0034}
  {993392326124fdc4d83c44ae53cdb333a2ca2c9911da027d164621951aafc62a}
\begin{o001statement}
Misalkan $S$ suatu fungsi signum (atau tanda) pada $\R$; artinya, misalkan $S = 2H - 1$, dengan $H$ fungsi Heaviside.
Misalkan pula $g$ suatu fungsi pada $\R$ sedemikian sehingga $g\colon x \mapsto -x - 3$ untuk $x < 0$ dan $g\colon x \mapsto x + 5$ untuk $x \ge 0$.
   \begin{enumerate}[label=\rm{(\alph*)}]
     \item Carilah konstanta $\alpha$ dan $\beta$ sedemikian sehingga ${\tilde g}' = \alpha \widetilde S + \beta \delta$.
     \item Carilah suatu fungsi $a$ pada $\R$ sedemikian sehingga ${\tilde a}' = \wt S$.
   \end{enumerate}
\end{o001statement}
\begin{o001answer}
$\alpha=1$, $\beta=8$, dan salah satu pilihan ialah $a(x)=\lvert x\rvert$.
Lebih umum, $a(x)=\lvert x\rvert+C$ untuk sembarang konstanta~$C$.
\end{o001answer}
\begin{o001proof}
Pada masing-masing setengah garis, turunan klasik $g$ sama dengan $-1$ untuk
$x<0$ dan $1$ untuk $x>0$, yaitu sama hampir di mana-mana dengan~$S$.
Sementara itu, loncatan $g$ di titik asal adalah
\[
  g(0+)-g(0-)=5-(-3)=8.
\]
Rumus turunan distribusional bagi fungsi mulus-sepotong karena itu memberi
\[
  {\tilde g}'=\widetilde S+8\delta,
\]
sehingga $\alpha=1$ dan $\beta=8$.

Untuk bagian kedua, ambil $a(x)=\lvert x\rvert$. Fungsi ini kontinu di nol,
jadi tidak menghasilkan suku delta, dan turunannya sama dengan $-1$ di
$(-\infty,0)$ serta $1$ di $(0,\infty)$. Maka
${\tilde a}'=\widetilde S$. Penambahan konstanta tidak mengubah turunan
distribusional.
\end{o001proof}
\end{o001solution}

% O001-SOLUTION-ID: O001-FAOA-2015-CH10-EX-002-SOLUTION
% SOURCE-EXERCISE-ID: FAOA-2015-CH10-NODE-0035
% STATEMENT-TARGET-SHA256: f48002bdb810711a169089b4f615782e5af6798ecf6f38a7be8141ef70a103e5
\begin{o001solution}
  {O001-FAOA-2015-CH10-EX-002-SOLUTION}
  {FAOA-2015-CH10-NODE-0035}
  {f48002bdb810711a169089b4f615782e5af6798ecf6f38a7be8141ef70a103e5}
\begin{o001statement}
Misalkan $f(x) = \left\{
                           \begin{array}{ll}
                             1 - \abs x, & \hbox{jika $\abs x \le 1$;} \\
                                 0,      & \hbox{selain itu.}
                           \end{array}
                         \right.$ Dengan menghitung kedua ruas persamaan~\eqref{deriv_reg_distr} secara terpisah, tunjukkan bahwa persamaan itu
benar secara klasik untuk setiap fungsi uji yang tumpuannya termuat dalam interval~$[-1,1]$.
\end{o001statement}
\begin{o001answer}
Kedua ruas bernilai
\[
  \int_{-1}^{0}\phi(x)\,dx-\int_{0}^{1}\phi(x)\,dx .
\]
\end{o001answer}
\begin{o001proof}
Turunan klasik $f$ ada hampir di mana-mana dan diberikan oleh
\[
 f'(x)=
 \begin{cases}
  1,&-1<x<0,\\
 -1,&0<x<1,\\
  0,&\text{selain itu}.
 \end{cases}
\]
Oleh sebab itu,
\[
 L_{f'}(\phi)=\int_{-1}^{0}\phi(x)\,dx-
               \int_{0}^{1}\phi(x)\,dx .
\]
Di sisi lain, dengan memecah integral di nol dan mengintegralkan per bagian,
\begin{align*}
 (L_f)'(\phi)
 &=-\int_{-1}^{0}(1+x)\phi'(x)\,dx
   -\int_{0}^{1}(1-x)\phi'(x)\,dx\\
 &=\left[-\phi(0)+\int_{-1}^{0}\phi(x)\,dx\right]
   +\left[\phi(0)-\int_{0}^{1}\phi(x)\,dx\right]\\
 &=\int_{-1}^{0}\phi(x)\,dx-\int_{0}^{1}\phi(x)\,dx.
\end{align*}
Suku batas di nol saling menghapus. Dengan demikian
$(L_f)'=L_{f'}$, sebagaimana dinyatakan oleh~\eqref{deriv_reg_distr} untuk
fungsi kontinu dan mulus-sepotong ini.
\end{o001proof}
\end{o001solution}

% O001-SOLUTION-ID: O001-FAOA-2015-CH10-EX-003-SOLUTION
% SOURCE-EXERCISE-ID: FAOA-2015-CH10-NODE-0036
% STATEMENT-TARGET-SHA256: b4a47ce3926ad911b776a8e31eb2efc713e5e925d51c317a06458160b0c7cbbc
\begin{o001solution}
  {O001-FAOA-2015-CH10-EX-003-SOLUTION}
  {FAOA-2015-CH10-NODE-0036}
  {b4a47ce3926ad911b776a8e31eb2efc713e5e925d51c317a06458160b0c7cbbc}
\begin{o001statement}
Misalkan $f(x) = \sbsb\chi{(0,1)}$ untuk $x \in~\R$ dan $\phi \in \fml D\bigl((-1,1)\bigr)$. Bandingkan nilai $L_{f\,'}(\phi)$ dan
$\bigl(L_f\bigr)'(\phi)$.
\end{o001statement}
\begin{o001answer}
$L_{f'}(\phi)=0$, sedangkan $(L_f)'(\phi)=\phi(0)$; jadi pada
$(-1,1)$ berlaku $(L_f)'=\delta_0$ dan keduanya tidak sama pada umumnya.
\end{o001answer}
\begin{o001proof}
Turunan klasik $f'$ sama dengan nol hampir di mana-mana, sehingga distribusi
reguler yang ditentukannya memenuhi $L_{f'}(\phi)=0$. Sebaliknya,
\[
 (L_f)'(\phi)=-L_f(\phi')=-\int_0^1\phi'(x)\,dx
 =\phi(0)-\phi(1)=\phi(0).
\]
Karena $\phi$ bertumpuan kompak di dalam $(-1,1)$, nilainya di~$1$ adalah
nol. Jadi turunan distribusional menangkap loncatan di~$0$, yang diabaikan
oleh turunan klasik hampir di mana-mana.
\end{o001proof}
\end{o001solution}

% O001-SOLUTION-ID: O001-FAOA-2015-CH10-EX-004-SOLUTION
% SOURCE-EXERCISE-ID: FAOA-2015-CH10-NODE-0037
% STATEMENT-TARGET-SHA256: 76fa760413de68845660f5155e23b736072da6a018d66a8e53d649fa3f39a8eb
\begin{o001solution}
  {O001-FAOA-2015-CH10-EX-004-SOLUTION}
  {FAOA-2015-CH10-NODE-0037}
  {76fa760413de68845660f5155e23b736072da6a018d66a8e53d649fa3f39a8eb}
\begin{o001statement}
Misalkan $f(x) = \left\{
                           \begin{array}{ll}
                             x, & \hbox{jika $0 \le x \le 1$;} \\
                             x + 2, & \hbox{jika $1 < x \le 2$;} \\
                             0, & \hbox{selain itu.}
                           \end{array}
                         \right.$ Bandingkan nilai $L_{f\,'}(\phi)$ dan $\bigl(L_f\bigr)'(\phi)$ untuk suatu fungsi uji~$\phi$.
\end{o001statement}
\begin{o001answer}
\[
 L_{f'}(\phi)=\int_0^2\phi(x)\,dx,
 \qquad
 (L_f)'(\phi)=\int_0^2\phi(x)\,dx+2\phi(1)-4\phi(2).
\]
\end{o001answer}
\begin{o001proof}
Di luar himpunan berukuran nol $\{0,1,2\}$, turunan klasik $f$ sama dengan
$1$ pada $(0,2)$ dan nol di luar interval tersebut. Maka
$L_{f'}(\phi)=\int_0^2\phi$. Fungsi $f$ tidak meloncat di~$0$, tetapi
loncatannya di~$1$ dan~$2$ berturut-turut adalah
\[
 f(1+)-f(1-)=3-1=2,
 \qquad
 f(2+)-f(2-)=0-4=-4.
\]
Turunan distribusional fungsi mulus-sepotong adalah distribusi reguler dari
turunan hampir-di-mana-mana ditambah besar setiap loncatan kali delta pada
titik loncatan. Karena itu
\[
 (L_f)'=L_{\sbsb\chi{(0,2)}}+2\delta_1-4\delta_2,
\]
yang memberikan rumus jawaban setelah diterapkan pada~$\phi$.
\end{o001proof}
\end{o001solution}

% O001-SOLUTION-ID: O001-FAOA-2015-CH10-EX-005-SOLUTION
% SOURCE-EXERCISE-ID: FAOA-2015-CH10-NODE-0038
% STATEMENT-TARGET-SHA256: f52569f7102a885e40cd395f3d064ee179df2aacdc336e54e891140188da963f
\begin{o001solution}
  {O001-FAOA-2015-CH10-EX-005-SOLUTION}
  {FAOA-2015-CH10-NODE-0038}
  {f52569f7102a885e40cd395f3d064ee179df2aacdc336e54e891140188da963f}
\begin{o001statement}
Misalkan $h(x) = \left\{
                           \begin{array}{ll}
                             2x, & \hbox{jika $x < -1$;} \\
                             1, & \hbox{jika $-1 \le x < 1$;} \\
                             0, & \hbox{jika $x \ge 1$.}
                           \end{array}
                         \right.$ Carilah suatu fungsi $f$ yang terintegralkan secara lokal sedemikian sehingga ${\tilde f}\,' = \tilde h + 3\delta_1 - \delta_{-1}$.
\end{o001statement}
\begin{o001answer}
Salah satu pilihan adalah
\[
 f(x)=
 \begin{cases}
 x^2,&x<-1,\\
 x+1,&-1\leq x<1,\\
 5,&x\geq1.
 \end{cases}
\]
Setiap solusi lain sama hampir di mana-mana dengan $f+C$ untuk suatu konstanta
global~$C$; pada tingkat distribusi, kebebasannya tepat satu konstanta global.
\end{o001answer}
\begin{o001proof}
Pada tiga interval terbuka yang bersangkutan, turunan klasik fungsi yang
diberikan berturut-turut adalah $2x$, $1$, dan~$0$, jadi bagian regulernya
adalah $\tilde h$. Di~$-1$, limit kanan dan kiri adalah
\[
 f(-1+)=0,
 \qquad
 f(-1-)=1,
\]
sehingga loncatannya $-1$. Di~$1$, limit kanan dan kiri adalah $5$ dan~$2$,
sehingga loncatannya~$3$. Oleh rumus turunan fungsi mulus-sepotong,
\[
 {\tilde f}\,'=\tilde h-\delta_{-1}+3\delta_1.
\]
Jika antiturunan pada ketiga interval diberi konstanta terpisah, syarat kedua
loncatan memaksa konstanta tersebut berbeda tepat seperti pada rumus di atas,
hingga penambahan satu konstanta yang sama pada semua bagian. Perubahan nilai
fungsi pada himpunan berukuran nol tidak mengubah distribusi reguler, sehingga
pernyataan keunikan untuk fungsi harus dipahami hampir di mana-mana.
\end{o001proof}
\end{o001solution}

% O001-SOLUTION-ID: O001-FAOA-2015-CH10-EX-006-SOLUTION
% SOURCE-EXERCISE-ID: FAOA-2015-CH10-NODE-0040
% STATEMENT-TARGET-SHA256: c64f5a0efcf4bce85f5462849dd1c213bdb4a6fde85b288a8891a0e99b5db78c
\begin{o001solution}
  {O001-FAOA-2015-CH10-EX-006-SOLUTION}
  {FAOA-2015-CH10-NODE-0040}
  {c64f5a0efcf4bce85f5462849dd1c213bdb4a6fde85b288a8891a0e99b5db78c}
\begin{o001statement}
Suatu \df{dipol} (dengan momen listrik 1 di titik asal pada $\R$) dapat dipandang sebagai ``limit'' ketika $\epsilon \sto 0^+$
 \index{dipol}%
dari suatu sistem $T_\epsilon$ yang terdiri atas dua muatan $-\frac1\epsilon$ dan $\frac1\epsilon$ yang masing-masing ditempatkan di $0$ dan $\epsilon$. Tunjukkan
bagaimana hal ini dapat mengarahkan kita untuk mendefinisikan dipol secara matematis sebagai $-\delta\,'$. \emph{Petunjuk.} Pandanglah $T_\epsilon$ sebagai
distribusi $\frac1\epsilon \delta_\epsilon - \frac1\epsilon \delta$.
\end{o001statement}
\begin{o001answer}
$T_\epsilon\to-\delta'$ dalam $\fml D^*(\R)$ ketika
$\epsilon\to0^+$.
\end{o001answer}
\begin{o001proof}
Untuk setiap $\phi\in\fml D(\R)$,
\[
 T_\epsilon(\phi)
 =\frac{\phi(\epsilon)-\phi(0)}{\epsilon}
 \longrightarrow \phi'(0).
\]
Menurut definisi turunan distribusi,
$\delta'(\phi)=-\delta(\phi')=-\phi'(0)$, sehingga
$(-\delta')(\phi)=\phi'(0)$. Jadi limit dua muatan tersebut, dalam arti
konvergensi titik demi titik pada setiap fungsi uji, tepat sama dengan
$-\delta'$.
\end{o001proof}
\end{o001solution}

% O001-SOLUTION-ID: O001-FAOA-2015-CH10-EX-007-SOLUTION
% SOURCE-EXERCISE-ID: FAOA-2015-CH10-NODE-0047
% STATEMENT-TARGET-SHA256: fd8775418583e8d49dd07a07086c23720b69fc2cd2a0946237eb620641240fcf
\begin{o001solution}
  {O001-FAOA-2015-CH10-EX-007-SOLUTION}
  {FAOA-2015-CH10-NODE-0047}
  {fd8775418583e8d49dd07a07086c23720b69fc2cd2a0946237eb620641240fcf}
\begin{o001statement}
Jika $\delta$ bukan suatu fungsi pada $\R$, lalu apa yang dimaksud orang ketika mereka menulis
   \[ \lim_{k \sto \infty} \frac{k}{\pi(1+k^2x^2)} = \delta(x)\quad? \]
Tunjukkan bahwa, jika ditafsirkan dengan tepat (sebagai suatu pernyataan tentang distribusi), ungkapan itu bahkan benar. \emph{Petunjuk.}
Jika $s_k(x) = k\,\pi^{-1}(1+k^2x^2)^{-1}$ dan $\phi$ suatu fungsi uji pada $\R$, maka
$\int_{-\infty}^\infty s_k\phi = \int_{-\infty}^\infty s_k\phi(0) + \int_{-\infty}^\infty s_k\eta$ dengan $\eta(x) = \phi(x) - \phi(0)$.
Tunjukkan bahwa $\int_{-\infty}^\infty s_k\eta \sto 0$ ketika $k \sto \infty$.
\end{o001statement}
\begin{o001answer}
Maknanya ialah $L_{s_k}\to\delta$ dalam $\fml D^*(\R)$, yaitu
\[
 \int_{\R}s_k(x)\phi(x)\,dx\longrightarrow\phi(0)
 \qquad\text{untuk setiap }\phi\in\fml D(\R).
\]
\end{o001answer}
\begin{o001proof}
Dengan substitusi $y=kx$ diperoleh
\[
 \int_{\R}s_k(x)\phi(x)\,dx
 =\frac1\pi\int_{\R}\frac{\phi(y/k)}{1+y^2}\,dy.
\]
Untuk setiap $y$, $\phi(y/k)\to\phi(0)$, sedangkan integrannya secara mutlak
dibatasi oleh
$\|\phi\|_\infty\big/\bigl(\pi(1+y^2)\bigr)$, sebuah fungsi yang
terintegralkan. Teorema konvergensi terdominasi memberi
\[
 \lim_{k\to\infty}\int_{\R}s_k\phi
 =\frac{\phi(0)}\pi\int_{\R}\frac{dy}{1+y^2}
 =\phi(0)=\delta(\phi).
\]
Pernyataan tersebut dengan demikian merupakan pernyataan tentang limit
distribusi, bukan limit titik demi titik fungsi.
\end{o001proof}
\end{o001solution}

% O001-SOLUTION-ID: O001-FAOA-2015-CH10-EX-008-SOLUTION
% SOURCE-EXERCISE-ID: FAOA-2015-CH10-NODE-0048
% STATEMENT-TARGET-SHA256: debe9140cb2f7d0021122f265e6c1e892dbe064fb8578e7594839762ed156c72
\begin{o001solution}
  {O001-FAOA-2015-CH10-EX-008-SOLUTION}
  {FAOA-2015-CH10-NODE-0048}
  {debe9140cb2f7d0021122f265e6c1e892dbe064fb8578e7594839762ed156c72}
\begin{o001statement}
Misalkan $f_n(x) = \sin nx$ untuk $n \in \N$ dan $x \in \R$. Barisan $(f_n)$ tidak konvergen secara klasik;
tetapi barisan itu konvergen secara distribusional (ke $0$). Nyatakan pernyataan ini dengan tepat, lalu buktikan. \emph{Petunjuk.} Carilah antiturunannya.
\end{o001statement}
\begin{o001answer}
Untuk setiap $\phi\in\fml D(\R)$,
$\int_{\R}\sin(nx)\phi(x)\,dx\to0$; yakni $L_{f_n}\to0$ dalam
$\fml D^*(\R)$.
\end{o001answer}
\begin{o001proof}
Karena $-\cos(nx)/n$ adalah antiturunan $\sin(nx)$ dan $\phi$ bertumpuan
kompak, integrasi per bagian menghasilkan
\[
 \int_{\R}\sin(nx)\phi(x)\,dx
 =\frac1n\int_{\R}\cos(nx)\phi'(x)\,dx.
\]
Oleh karena itu,
\[
 \left|\int_{\R}\sin(nx)\phi(x)\,dx\right|
 \leq\frac{\|\phi'\|_{L^1}}{n}\longrightarrow0.
\]
Inilah tepat definisi konvergensi distribusional ke distribusi nol.
\end{o001proof}
\end{o001solution}

% O001-SOLUTION-ID: O001-FAOA-2015-CH10-EX-009-SOLUTION
% SOURCE-EXERCISE-ID: FAOA-2015-CH10-NODE-0050
% STATEMENT-TARGET-SHA256: ec002d6bdad3dd72b45e60fa31af1c8ef543b8b53ad21cee638924f22b6cfd48
\begin{o001solution}
  {O001-FAOA-2015-CH10-EX-009-SOLUTION}
  {FAOA-2015-CH10-NODE-0050}
  {ec002d6bdad3dd72b45e60fa31af1c8ef543b8b53ad21cee638924f22b6cfd48}
\begin{o001statement}
Berikan makna pada ungkapan
   \[ 2\pi\sum_{n = -\infty}^\infty \delta(x-2\pi n) =  1 +  2\sum_{n = 1}^\infty \cos nx\,, \]
dan tunjukkan bahwa, jika ditafsirkan dengan tepat, ungkapan itu benar. \emph{Petunjuk.} Siapa pun yang cukup bejat untuk percaya bahwa $\delta(x)$ \emph{berarti}
sesuatu kemungkinan akan menafsirkan $\delta(x-a)$ sebagai hal yang sama dengan $\delta_a(x)$. Anda dapat memulai proses
merasionalisasikannya dengan menghitung deret Fourier fungsi $g$ yang didefinisikan oleh $g(x) = \frac1{4\pi}x^2 - \frac12 x$ pada $[0,2\pi]$
dan kemudian diperluas secara periodik ke $\R$. Anda mungkin perlu mencari teorema tentang konvergensi titik demi titik dan seragam dari deret Fourier
serta tentang pendiferensialan suku demi suku deret semacam itu.
\end{o001statement}
\begin{o001answer}
Dalam $\fml D^*(\R)$ berlaku
\[
 2\pi\sum_{k\in\Z}\delta_{2\pi k}
 =\lim_{N\to\infty}\left(1+2\sum_{n=1}^{N}\cos(nx)\right).
\]
Kedua ruas yang diterapkan pada $\phi$ bernilai
$2\pi\sum_{k\in\Z}\phi(2\pi k)$.
\end{o001answer}
\begin{o001proof}
Ambil $\phi\in\fml D(\R)$ dan periodisasikan fungsi itu dengan
\[
 P\phi(t)=\sum_{k\in\Z}\phi(t+2\pi k).
\]
Pada setiap titik hanya berhingga banyak suku yang tidak nol secara lokal,
sehingga $P\phi$ mulus dan $2\pi$-periodik. Koefisien Fourier-nya adalah
\[
 c_n=\frac1{2\pi}\int_0^{2\pi}P\phi(t)e^{-int}\,dt
 =\frac1{2\pi}\int_{\R}\phi(t)e^{-int}\,dt.
\]
Integrasi per bagian berulang menunjukkan bahwa $(c_n)$ menurun lebih cepat
daripada setiap pangkat $|n|^{-1}$. Jadi deret Fourier tersebut konvergen
mutlak dan seragam. Dengan mengevaluasinya di~$0$, kita memperoleh rumus
penjumlahan Poisson
\[
 2\pi\sum_{k\in\Z}\phi(2\pi k)
 =\sum_{n\in\Z}\int_{\R}\phi(t)e^{-int}\,dt.
\]
Penggantian $n$ oleh $-n$ tidak mengubah jumlah atas semua bilangan bulat,
sehingga ruas kanan juga sama dengan
\[
 \lim_{N\to\infty}\int_{\R}
 \left(\sum_{n=-N}^{N}e^{int}\right)\phi(t)\,dt
 =\lim_{N\to\infty}\int_{\R}
 \left(1+2\sum_{n=1}^{N}\cos(nt)\right)\phi(t)\,dt.
\]
Sementara itu, ruas kiri adalah penerapan distribusi
$2\pi\sum_{k\in\Z}\delta_{2\pi k}$ pada~$\phi$; jumlah ini berhingga untuk
setiap fungsi uji. Identitas berlaku untuk setiap~$\phi$, sehingga berlaku
dalam arti distribusi.
\end{o001proof}
\end{o001solution}

% O001-SOLUTION-ID: O001-FAOA-2015-CH10-EX-010-SOLUTION
% SOURCE-EXERCISE-ID: FAOA-2015-CH10-NODE-0086
% STATEMENT-TARGET-SHA256: 0f28bc54f7a5f5bffbd03ec67bc3d3d57cd2e3851430c8c294375d7fe333d618
\begin{o001solution}
  {O001-FAOA-2015-CH10-EX-010-SOLUTION}
  {FAOA-2015-CH10-NODE-0086}
  {0f28bc54f7a5f5bffbd03ec67bc3d3d57cd2e3851430c8c294375d7fe333d618}
\begin{o001statement}
Buktikan proposisi sebelumnya dengan asumsi bahwa $u$ dan $v$ \emph{keduanya} memiliki tumpuan kompak. \emph{Petunjuk.}
Gunakan $(u \ast v) \ast (\phi \ast \psi)$, dengan $\phi$ dan $\psi$ fungsi-fungsi uji, sambil mengingat bahwa konvolusi fungsi
bersifat komutatif.
\end{o001statement}
\begin{o001answer}
Jika $u$ dan $v$ bertumpuan kompak, maka $u\ast v=v\ast u$.
\end{o001answer}
\begin{o001proof}
Karena $u$ dan $v$ bertumpuan kompak, $u\ast\psi$ dan $v\ast\phi$
merupakan fungsi uji untuk setiap $\phi,\psi\in\fml D(\R^n)$. Dengan definisi
konvolusi dua distribusi, sifat asosiatif konvolusi distribusi--fungsi-uji,
dan komutativitas konvolusi fungsi, diperoleh
\begin{align*}
 (u\ast v)\ast(\phi\ast\psi)
 &=u\ast\bigl(v\ast(\phi\ast\psi)\bigr)\\
 &=u\ast\bigl((v\ast\phi)\ast\psi\bigr)\\
 &=u\ast\bigl(\psi\ast(v\ast\phi)\bigr)\\
 &=(u\ast\psi)\ast(v\ast\phi).
\end{align*}
Di sisi lain, karena $\phi\ast\psi=\psi\ast\phi$, perhitungan yang sama
dengan $u,v$ dan $\phi,\psi$ dipertukarkan memberi
\[
 (v\ast u)\ast(\phi\ast\psi)
 =(v\ast\phi)\ast(u\ast\psi).
\]
Dua ruas terakhir sama karena keduanya merupakan konvolusi fungsi uji.

Sekarang tetapkan $w=u\ast v-v\ast u$. Ambil sembarang
$\eta\in\fml D(\R^n)$ dan suatu aproksimasi identitas
$(\rho_\epsilon)$ dalam $\fml D(\R^n)$. Hasil di atas dengan
$\phi=\eta$ dan $\psi=\rho_\epsilon$ menunjukkan
\[
 w\ast(\eta\ast\rho_\epsilon)=0.
\]
Karena $\eta\ast\rho_\epsilon\to\eta$ dalam topologi fungsi uji dan distribusi
kontinu pada topologi itu, kita memperoleh $w\ast\eta=0$. Hal ini berlaku
untuk setiap~$\eta$. Proposisi ketunggalan tepat sebelum definisi tumpuan
kemudian memberi $w=0$, yaitu $u\ast v=v\ast u$.
\end{o001proof}
\end{o001solution}

% O001-SOLUTION-ID: O001-FAOA-2015-CH10-EX-011-SOLUTION
% SOURCE-EXERCISE-ID: FAOA-2015-CH10-NODE-0099
% STATEMENT-TARGET-SHA256: d7c771909a0b07a97813b77756de75190669a00de4df0212ffbcdc10f4c996f3
\begin{o001solution}
  {O001-FAOA-2015-CH10-EX-011-SOLUTION}
  {FAOA-2015-CH10-NODE-0099}
  {d7c771909a0b07a97813b77756de75190669a00de4df0212ffbcdc10f4c996f3}
\begin{o001statement}
Carilah suatu solusi distribusional bagi persamaan diferensial $x \dfrac{du}{dx} + u = 0$ pada~$\R$.
\end{o001statement}
\begin{o001answer}
Distribusi delta Dirac $u=\delta$ merupakan suatu solusi.
\end{o001answer}
\begin{o001proof}
Untuk setiap fungsi uji~$\phi$,
\[
 (x\delta')(\phi)=\delta'(x\phi)
 =-(x\phi)'(0)=-\phi(0)=-\delta(\phi).
\]
Jadi $x\delta'=-\delta$, dan akibatnya
\[
 x\frac{d\delta}{dx}+\delta=x\delta'+\delta=0
\]
dalam $\fml D^*(\R)$.
\end{o001proof}
\end{o001solution}
